\documentclass[11pt]{article}
\usepackage[a4paper,margin=26mm]{geometry}
\usepackage{amsmath,amssymb,amsthm,enumitem,hyperref}
\newtheorem{theorem}{Theorem}
\newtheorem{lemma}[theorem]{Lemma}
\newtheorem{remark}[theorem]{Remark}
\newcommand{\Int}{\operatorname{Int}}
\newcommand{\htS}{h_S}
\newcommand{\htT}{h_T}
\newcommand{\htQ}{h_Q}
\newcommand{\pre}{\sqsubseteq}
\newcommand{\lex}{\le_{\mathrm{lex}}}
\newcommand{\fwd}{\preccurlyeq_{\mathrm{fwd}}}
\newcommand{\len}[1]{\lvert #1\rvert}

\title{A conditional geometric repair of the Lemma 27 \(Q\)-construction}
\author{Verification note for Todorcevic--Tyros,
\emph{A Dual Ramsey Theorem for Trees}}
\date{October 2026}
\begin{document}
\maketitle

\paragraph{Scope.}
This note states a \emph{corrected and explicitly conditional}
version of the tree-support component of Lemma~27
(arXiv:2207.14599v1).
It is not an erratum for the entire paper and does not
prove its Ramsey conclusion. In particular, neither the
mixed-cone word \(g_w\), its smoothness/colouring pullback,
nor the original recursive bounds are supplied here.
The printed length-first reverse-lexicographic order and
unrestricted leaf replacement have separately documented
counterexamples in
\texttt{notes/lemma27\_Q\_counterexample.tex}.
The development below uses the globally \emph{forward}
length-first auxiliary order consistently.

\section*{Data and precise hypothesis}

Let \(b>0\), and let \(S\) be a finite semi-complete skew
source support with its last interior vertex \(t_0\)
in length-first forward auxiliary order.
The original starred source word has variable roots
\(\Int(S)\).
Let \(T\) be a finite \(k\)-complete skew support under
the same forward convention, with \(S\subseteq T\) and
\(t_0\in T\).
Write \(m=\htT(t_0)\), and assume
\begin{equation}\label{eq:early}
 m+1<k.
\end{equation}
The repaired signature has a literal marker set \(R\)
comprising (a) the first exit boundary of each exceptional
terminal \(t\in S\) not extending \(t_0\), and (b) each
ambient child \(t_0^\frown i\) for \(i<b\).
Every marker has a unique minimal \(T\)-frontier \(f\)
outside the \emph{inclusive} forward cut at \(t_0\).
The marker-to-frontier assignment is injective.
Enumerate \emph{all} such minimal frontiers in increasing
forward auxiliary order, and let \(D_2\) be the finite
subset of indices assigned to markers in \(R\).

For each \(i\in D_2\), choose a local marked point
\(s_i\in b^{<k-(m+1)}\). Let \(f_i\) be its selected
\(T\)-frontier, and let
\[
 P_i=I_T\bigl(I_T^{-1}(f_i)^\frown s_i\bigr),
\]
where \(I_T:b^{<k}\to T\) is the canonical
support-address map for the \emph{forward} complete
skew support.
The shortened common tail space makes every \(P_i\)
well-defined even when \(\htT(f_i)=m+1\), and
\begin{equation}\label{eq:projectrank}
 f_i\pre P_i,\qquad
 \htT(P_i)=\htT(f_i)+\len{s_i}.
\end{equation}
The actual reconstructed support is
\[
 Q=\Int(S)\cup\{P_i:i\in D_2\}.
\]

\paragraph{Starred \(D_2\)-order hypothesis.}
The only numerical consequence of Definition~20 used
by the geometric argument is
\begin{equation}\label{eq:D2}
  i<j,\quad i,j\in D_2
       \quad\Longrightarrow\quad
       \len{s_i}\le\len{s_j}.
\end{equation}
This follows from a genuine vector \(1\)-complete
singleton family indexed by the \emph{selected} \(D_2\)
coordinates in their induced order. It is false to
assume \eqref{eq:D2} for an arbitrary element of the
unstarred mixed product. The Lean development uses an
explicit order-reflecting enumeration/certificate and
does not silently identify it with every printed
\(W_*\)-input.

\section*{The matching of the two intrinsic rank bands}

\begin{lemma}[Exact leaf/frontier rank matching]
\label{lem:bands}
For every actual marked projection \(P_i\) one of the
following two alternatives holds:
\[
\begin{aligned}
 \bigl(\htQ(P_i),\htT(f_i)\bigr)
      &=(\htS(t_0),\htT(t_0)),\\
 \text{or}\qquad
 \bigl(\htQ(P_i),\htT(f_i)\bigr)
      &=(\htS(t_0)+1,\htT(t_0)+1).
\end{aligned}
\]
The alternative is independent of the local point \(s_i\).
\end{lemma}

\begin{proof}
All new projected \(P_i\) are terminal in \(Q\),
and the actual old interior is preserved:
\(\Int(Q)=\Int(S)\). Consequently \(\htQ(P_i)\)
is the number of strict old-interior predecessors of
its literal marker.

For an exceptional terminal \(t\in S\), the first-exit
boundary has exactly the original interior predecessors
of \(t\); thus \(\htQ(P_i)=\htS(t)\).
A minimal outside-cut source terminal has height
\(\htS(t_0)\) or \(\htS(t_0)+1\).

Its selected \(T\)-frontier \(f_i\) lies
\emph{on the original path to \(t\)}:
both \(t\) and \(f_i\) extend the literal first boundary,
so their meet is still outside the inclusive cut.
Meet-closure of \(T\) and minimality of \(f_i\)
force \(f_i\pre t\).

If \(\htS(t)=\htS(t_0)+1\), choose the genuine
source-support ancestor \(u\pre t\) of height
\(\htS(t_0)\). The two same-height interior points
\(u,t_0\) are in incomparable cones and
\(u\lex t_0\). This lexicographic relation persists
to \(f_i\), and the existing strict boundary-length
argument shows that the first exit is longer than \(t_0\).
Were \(\htT(f_i)=\htT(t_0)\), the \(T\) skew
same-rank clause would put \(f_i\fwd t_0\),
contradicting its being outside the inclusive cut.
The general frontier bound
\(\htT(f_i)\le\htT(t_0)+1\) gives the high alternative.

If \(\htS(t)=\htS(t_0)\), source skewness gives
\(t_0\lex t\). Take the \(T\)-support ancestor \(a\pre t\)
of height \(\htT(t_0)\). It is incomparable with \(t_0\)
and lexicographically to its right, hence outside
the inclusive cut by \(T\) skewness. The frontier
\(f_i\) must prefix \(a\), or \(a\) would violate
its outside-cut minimality. Thus
\(\htT(f_i)\le\htT(a)=\htT(t_0)\), and the standard
outside-cut frontier lower bound gives equality.

Finally, if the literal marker is the direct child
\(t_0^\frown j\), then \(t_0\) is itself a strict
predecessor in \(Q\), giving
\(\htQ(P_i)=\htS(t_0)+1\).
Its selected frontier extends \(t_0^\frown j\),
hence has \(\htT(f_i)>\htT(t_0)\).
The same frontier upper bound gives
\(\htT(f_i)=\htT(t_0)+1\).
These are all the literal markers.
\end{proof}

\begin{remark}
Lemma~\ref{lem:bands} is stronger than the previously
known band bounds on \(Q\)-heights and \(T\)-frontier
heights separately. It proves that the \emph{same}
alternative occurs for each selected marker.
Therefore
\[
\begin{gathered}
 \htQ(P_i)<\htQ(P_j)
   \iff \htT(f_i)<\htT(f_j),\\
 \htQ(P_i)=\htQ(P_j)
   \iff \htT(f_i)=\htT(f_j).
\end{gathered}
\]
\end{remark}

\section*{Conditional semi-completeness}

\begin{theorem}[Corrected \(Q\)-support component]
\label{thm:Q}
Under the data above, including the explicit
\(D_2\)-length order \eqref{eq:D2},
the actual reconstructed support \(Q\) is a
semi-complete skew tree under the
length-first \emph{forward} auxiliary order.
\end{theorem}

\begin{proof}
The root, exact old interior, ambient meet closure,
\(0\)-or-\(b\) immediate successor counts, all
directions below the maximal cut, and final skew
branching witness are independent verified geometric
facts about this corrected \(Q\) support.
It therefore suffices to verify the two skew
rank/ambient-length conditions.

For two old interior vertices, both clauses transfer
from \(S\), since their intrinsic heights in \(Q\)
are unchanged. For an old interior vertex and a
marked projection, the equal-rank lexicographic
orientation and strict-rank length inequality
have been proved using the literal exceptional and
cut-child marker cases. A new marked leaf can never
have smaller \(Q\)-height than an old interior vertex.

It remains to compare two new projected points.
The genuine forward-sorted frontier list has
nondecreasing \(T\)-ranks by index.
If \(\htQ(P_i)=\htQ(P_j)\), Lemma~\ref{lem:bands}
gives \(\htT(f_i)=\htT(f_j)\).
The selected frontier cones are incomparable;
therefore \(P_i\lex P_j\) forces \(i\le j\).
The case \(i=j\) is trivial. If \(i<j\),
hypothesis \eqref{eq:D2} and the verified
monotonicity of canonical projections in the
sorted frontier order give
\(\len{P_i}\le\len{P_j}\).
This is the new/new skew same-height clause.

If \(\htQ(P_i)<\htQ(P_j)\), then
\(\htT(f_i)<\htT(f_j)\) by
Lemma~\ref{lem:bands}. Frontier-rank
monotonicity excludes \(j\le i\), hence \(i<j\).
Using \eqref{eq:projectrank} and \eqref{eq:D2},
\[
 \htT(P_i)=\htT(f_i)+\len{s_i}
         <\htT(f_j)+\len{s_j}
         =\htT(P_j).
\]
The skew strict-height clause in the enclosing
complete support \(T\) now yields
\(\len{P_i}<\len{P_j}\).
This is the remaining new/new condition.
\end{proof}

\section*{What is still not proved}

\begin{enumerate}[label=(\roman*)]
\item The source-facing Lean certificate of
selected \(D_2\) vector \(1\)-completeness
must be constructed from precisely the actual
members of the printed \(W_*\) domain.
This note does not assert the same length property
for arbitrary mixed-product inputs.
\item The paper's original local cone height
and recursive numerical bounds must be reconciled
with the shortened common tail space in
\eqref{eq:early}, or an extra complete level
must be supplied and the iteration adjusted.
\item The particular word \(g_w\) must be constructed,
including root-aligned substitution, typed
three-way local symbol decoding, and
consistency of its variable fibres.
\item Smoothness of the colouring pullback
and classwise compatibility must be proved
before applying mixed-product Ramsey results.
A separate checked example shows that an
arbitrary letter-at-marked-point colouring
need not be smooth.
\item Lemma~28, Theorem~25, and the full dual
Ramsey theorem remain beyond this note.
\end{enumerate}

\paragraph{Lean status and source policy.}
The statements in this note are mirrored by the
modules \texttt{ForwardQFrontierBandMatch} and
\texttt{ForwardQCertifiedSemiComplete}, with supporting
source-facing modules for each marker family.
Proof status must be read from the corresponding
current GitHub Actions runs and transitive axiom audit.
Until those complete, the statements above are
\emph{proposed Lean proofs awaiting CI}.
The original paper's manuscript is not modified;
proposed corrections and verification markers
are maintained in \texttt{audit.tex}.

\end{document}
